A common term used to access the sampling quality is the so-called \emph{ergodic principle}. Referring to \cref{eq:general_sampling_2}, the thermodynamic averaging of the property of a material is supposed to be exact when the MD sampling time $t$ or the number of MC steps $N$ goes to infinity:
\begin{equation}
    \label{eq:ergody_condition}
    \langle\mathcal{A}\rangle_T=\lim_{t\to\infty}\int_0^t\mathcal{A}(\bm{R}(t))\mbox{d}t =\int \mathcal{A}(\bm{R}) P(\bm{R})\mbox{d}\bm{R}.
\end{equation}
However, ergodicity can only be met if the underlying PES being sampled is relatively flat with shallow minima separated by low energy barriers. In realistic materials systems, very often, one sees deep local minima/global minimum that are separated from the main part of the energy landscape with high energy barriers [see, for example, \cref{fig:energy_landscape_examples}(b)]. On one hand, the low energy of these configurations means they contribute significantly to the Boltzmann weighting, and are important `landmarks' of the PES. On the other hand, the high-energy barriers translate to a low probability or even the impossibility of visiting these deep minima in an equilibrium MC/MD simulation. Visiting these deep minima thus becomes the `rare events' in equilibrium samplings, and overcoming this limit is the key driving force behind the development of the enhanced sampling methods in atomistic modelling, which we will review some of the major approaches that have been developed in the field.